\textbf{Acquired actor coefficients, opportunity contact and expenditure have distinct roles in the tested futures.}
\textbf{a,} Representative compact history h2 from the second opportunity round. The four body fields are archived states 0, 256 and 512 of acquisition, and state 768 of the B/retain continuation. Each is shown in the same 24 by 24 cell window, x6--29 and y2--25, in the original 32 by 32 periodic coordinates with y increasing upward. The field is neither translated with the body nor normalized by its own maximum. Density has a fixed linear range 0--1 tissue units per cell. Native outlines locate A (x14--18) and B (x21--25), extending through the full domain in y; they do not depict resource stock or uptake. Only A can supply during acquisition. The B target becomes available at action 513, subsequently for 16 actions in every 64. B already contains body mass during acquisition, partly from initial tails; the lower trace shows its rise and subsequent decline before B is available. Contact is current tissue mass in the entire strip, not a count of visits. The retain and reset contact traces share the acquisition prefix and differ only after the common receiver at 512. This is not a first-discovery or actuator-inaccessibility example. Cropped images do not define the spatial support of any total.

\textbf{b,} Full-domain tissue M, stored energy E, and cumulative uptake U and realized motor payment for the same compact h2 B future. The latter two accumulate from actions 513--768, not over acquisition. Both branches begin with E=3.220887 EU; retain/reset end with M=3.765989/1.232265 and E=0.344504/0.002881 EU. Uptake is only 0.349571/0.335341 EU, with motor payment 1.558061/2.362788 EU. Retention therefore accompanies a large tissue difference while intake remains small and the initial reserve is largely depleted. These curves do not identify a unique mediator of the endpoint difference. Each quantity has its labelled absolute scale; cumulative payment is a whole-body total. Lines are editable, simplified polylines through archived per-action readouts, rather than new forward calculations.

\textbf{c,} Signed endpoint and integrated-payment differences for all four genuine acquisition histories (two compact, two annular) crossed with future availability at A or B. The second physical construction was selected adaptively. These are eight paired conditions, not eight independent acquisition histories or 24 independent replicates. Each reached state is tested with retained, zeroed or shared actor coefficients under the same future site and random tape. All actor increments are disabled during these tests; transport, critic, food, score trace and m/z updates continue. There is no old T controller: run/brake modulates the transported polarity readout and does not stop the physical dynamics. Teal denotes retain minus reset for all four displayed quantities; orange denotes shared minus reset for energy. The first three columns share an EU scale, whereas tissue differences use a separately labelled scale. Tiny negative tissue effects are also identified by their signs. All eight retained-state energy effects are positive and motor payments lower, but compact h2/A has less uptake and three A pairs have negative total-mass effects despite more final A contact. Shared coefficients also give positive energy effects against reset in all eight pairs; the retained spatial field does not uniformly outperform them. Shared coefficients are externally calibrated to match the initial mass-weighted mean run probability while preserving the mean dry/fed contrast. This matches neither subsequent actions nor costs; local food continues to enter execution.

\textbf{d,} Separate earlier experiment: two acquisition schedules provide four within-receiver retain/reset comparisons over actions 513--640. Retained-state energy effects are negative under continuous access and positive under intermittent access for both histories. These receivers differ in physical state and stock across histories. For the intermittent-acquisition receiver, the retained/reset intermittent future ends with M=22.419509/22.251675, E=23.407757/21.760639 EU and U=12.372958/12.753755 EU. The continuous-history intermittent comparison is still energy-negative at 576 and positive at 640; no first sign-change time is inferred. Orange marks add the shared-coefficient diagnostic on the intermittent-history receiver, retaining the continuous-negative/intermittent-positive pattern. All earlier tested futures lose tissue. Initial mean probability matching already permits different realized actions at 513. The second-round eight positive energy effects do not erase these earlier conditional counterexamples. No panel establishes general adaptation, spatial mediation, or the value of continued actor learning. All fields and readouts reuse completed archives: second round 8192 updates; earlier transfer 2048; shared diagnostic 256 new updates with four earlier futures reused. Figure preparation adds zero model or energy updates.
